% ============ 第3章 §3.2 若干交换Banach代数的Gelfand表示(原书 pp.226-237) ============
\section{若干交换 Banach 代数的 Gelfand 表示}\label{sec:3.2}
\index{Fourier 级数}
\index{Fourier 系数}
\index{测度代数}
\index{Fourier--Stieltjes 变换}
\index{乘子!(单位圆周上的函数的)}
\index{乘子!乘子算子(单位圆周上的函数的)}
\index{Calderón--Zygmund 定理}
\index{Katznelson 问题}
\index{Parreau 定理}
\index{Wiener--Pitt 现象}
\index{Williamson 定理}


本节作为 Fourier 分析中出现的交换 Banach 代数的例子,考察具有绝对收敛 Fourier 级数的连续函数空间,Fourier 变换的空间等,并叙述它们的极大理想空间.另外,Riemann-Lebesgue 引理等初等 Fourier 分析的知识视为已知.

\subsection{\texorpdfstring{关于 $A(\mathbf{T})$}{关于 A(T)}}\label{sec:3.2.1}

\begin{definition}\label{def:3.2.1}
定义
\begin{equation*}
A(\mathbf{T}) = \left\{f \in C(\mathbf{T}):\ \|f\|_{A(\mathbf{T})} = \sum_{n=-\infty}^{\infty}|\hat{f}(n)| < \infty\right\}.
\end{equation*}
其中,
\begin{equation*}
\hat{f}(n) = \frac{1}{2\pi}\int_{0}^{2\pi}f(x)e^{-inx}\dd x
\end{equation*}
(一般地,对可积函数 $f$,称 $\hat{f}(n)$ 为 $f$ 的 Fourier 系数,称形式级数 $\sum_{n \in \mathbf{Z}}\hat{f}(n)e^{inx}$ 为 $f$ 的 Fourier 级数).又,和与数乘定义为
\begin{equation*}
(\alpha f + \beta g)(x) = \alpha f(x) + \beta g(x) \quad (\alpha, \beta \in \mathbf{C},\ f, g \in A(\mathbf{T})),
\end{equation*}
积定义为
\begin{equation*}
(fg)(x) = f(x)g(x) \quad (f, g \in A(\mathbf{T})),
\end{equation*}
则 $A(\mathbf{T})$ 成为环.进一步,$\|\cdot\|_{A(\mathbf{T})}$ 满足范数的性质,且关于这个范数完备.又,$\|fg\|_{A(\mathbf{T})} \leq \|f\|_{A(\mathbf{T})}\|g\|_{A(\mathbf{T})}$ 成立,故由具有绝对收敛 Fourier 级数的连续函数构成的空间 $A(\mathbf{T})$ 是以常数函数 $1$ 为单位元的交换 Banach 代数.
\end{definition}

\begin{theorem}\label{thm:3.2.2}
$A(\mathbf{T})$ 的极大理想空间 $\mathfrak{M}(A(\mathbf{T}))$ 与 $\mathbf{T}$ 等同.Gelfand 变换是恒等映射.又,$A(\mathbf{T})$ 半单.
\end{theorem}

\begin{proof}
由 $\mathfrak{M}(A(\mathbf{T})) = \Delta(A(\mathbf{T}))$,考察 $\chi \in \Delta(A(\mathbf{T}))$.因 $e^{i\theta} \in A(\mathbf{T})$,
\begin{equation*}
\chi(e^{i\theta})\chi(e^{-i\theta}) = \chi(e^{i\theta}e^{-i\theta}) = \chi(1) = 1
\end{equation*}
成立.因此 $\chi(e^{i\theta})\chi(e^{-i\theta}) = 1$.进一步,$|\chi(e^{i\theta})| \leq \|e^{i\theta}\|_{A(\mathbf{T})} = 1$,故 $|\chi(e^{i\theta})| = 1$.从而可令 $\chi(e^{i\theta}) = e^{i\xi}$($\xi \in [0, 2\pi)$).于是,对 $f \in A(\mathbf{T})$,$f(\theta) = \sum_{n \in \mathbf{Z}}\hat{f}(n)e^{in\theta}$,令 $f_{N}(\theta) = \sum_{|n| \leq N}\hat{f}(n)e^{in\theta}$,由 $\|f_{N} - f\|_{A(\mathbf{T})} \to 0$($N \to \infty$),
\begin{equation*}
|\chi(f) - \chi(f_{N})| = |\chi(f - f_{N})| \leq \|f_{N} - f\|_{A(\mathbf{T})} \to 0 \quad (N \to \infty)
\end{equation*}
成立.因此
\begin{equation*}
\chi(f) = \sum_{n \in \mathbf{Z}}\hat{f}(n)e^{in\xi}.
\end{equation*}
即,对 $\chi \in \Delta(A(\mathbf{T}))$ 对应 $\xi \in [0, 2\pi)$.反之,对 $\xi \in [0, 2\pi)$ 考虑 $\chi_{\xi}:\ f \mapsto f(\xi)$,则 $\chi_{\xi} \in \Delta(A(\mathbf{T}))$.因此,由 $\xi \in [0, 2\pi)$ 到 $\chi_{\xi} \in \Delta(A(\mathbf{T}))$ 的对应是满射.由 $\chi_{\xi}(f) = \chi_{\xi'}(f)$($f \in A(\mathbf{T})$)得 $e^{i\xi} = e^{i\xi'}$($\xi, \xi' \in [0, 2\pi)$),故 $\xi = \xi'$,是单射.又,当 $\lim_{n \to \infty}\xi_{n} = \xi$ 时,
\begin{equation*}
\lim_{n \to \infty}\chi_{\xi_{n}}(f) = \lim_{n \to \infty}f(\xi_{n}) = f(\xi)
\end{equation*}
故
\begin{equation*}
\lim_{n \to \infty}\chi_{\xi_{n}}(f) = \chi_{\xi}(f)
\end{equation*}
成立.因 $\mathbf{T}$ 是紧 Hausdorff 空间,$\Delta(A(\mathbf{T}))$ 与 $\mathbf{T}$ 同胚.因此 Gelfand 变换是恒等映射.由此,$A(\mathbf{T})$ 半单.
\end{proof}

\begin{note}\label{note:3.2.3}
$n \in \mathbf{N}$,令 $\mathbf{T}^{n} = \{\theta = (\theta_{1}, \ldots, \theta_{n}):\ \theta_{j} \in [0, 2\pi)\ (j = 1, 2, \ldots, n)\}$.定义~\ref{def:3.2.1} 的多维版本定义为
\begin{equation*}
A(\mathbf{T}^{n}) = \left\{f \in C(\mathbf{T}^{n}):\ f(\theta) = \sum_{m = (m_{1}, \ldots, m_{n}) \in \mathbf{Z}^{n}}\hat{f}(m)e^{im \cdot \theta},\ \|f\|_{A(\mathbf{T}^{n})} = \sum_{m \in \mathbf{Z}^{n}}|\hat{f}(m)| < \infty\right\}.
\end{equation*}
其中 $m \cdot \theta = \sum_{j=1}^{n}m_{j}\theta_{j}$,
\begin{equation*}
\hat{f}(m) = \frac{1}{(2\pi)^{n}}\int_{0}^{2\pi}\int_{0}^{2\pi}\cdots\int_{0}^{2\pi}f(\theta_{1}, \ldots, \theta_{n})e^{-im \cdot \theta}\dd\theta.
\end{equation*}
此时,对 $A(\mathbf{T}^{n})$ 也有与定理~\ref{thm:3.2.2} 相同的结论.
\end{note}

\subsection{\texorpdfstring{Fourier 变换构成的交换 Banach 代数 $A(\mathbb{R}^n)$}{Fourier 变换构成的交换 Banach 代数 A(R^n)}}\label{sec:3.2.2}

\begin{definition}\label{def:3.2.4}
$n$ 维 Euclid 空间 $\mathbf{R}^{n}$ 上的 Lebesgue 可积函数空间定义为
\begin{equation*}
L^{1}(\mathbf{R}^{n}) = \left\{f:\ \|f\|_{1} = \frac{1}{(2\pi)^{n}}\int_{\mathbf{R}^{n}}|f(x)|\dd x < \infty\right\}.
\end{equation*}
进一步,由
\begin{equation*}
\hat{f}(\xi) = \frac{1}{(2\pi)^{n}}\int_{\mathbf{R}^{n}}f(x)e^{-i\xi \cdot x}\dd x
\end{equation*}
定义 $f$ 的 Fourier 变换.其中 $\xi \cdot x$ 是内积.
\end{definition}

对 $f, g \in L^{1}(\mathbf{R}^{n})$,$\alpha, \beta \in \mathbf{C}$,由 $(\alpha f + \beta g)(x) = \alpha f(x) + \beta g(x)$,以及 $f * g(x) = \frac{1}{(2\pi)^{n}}\int_{\mathbf{R}^{n}}f(x - y)g(y)\dd y$,$L^{1}(\mathbf{R}^{n})$ 关于和,数乘,积成为环.又,$L^{1}(\mathbf{R}^{n})$ 以 $\|\cdot\|_{1}$ 为范数成为 Banach 空间,且满足 $\|f * g\|_{1} \leq \|f\|_{1}\|g\|_{1}$.综上,$L^{1}(\mathbf{R}^{n})$ 成为交换 Banach 代数.

又,对测度空间 $(X, \mu)$,用 $L^{\infty}(X)$ 表示以 $\|f\|_{L^{\infty}} = \inf\{M:\ |f(x)| < M\ (\mu\text{-a.e.})\}$ 为范数的 Banach 空间.以后,$L^{1}(\mathbf{R})^{*} = L^{\infty}(\mathbf{R})$,$L^{1}(\mathbf{T})^{*} = L^{\infty}(\mathbf{T})$ 等视为已知.

\begin{theorem}\label{thm:3.2.5}
下列事实成立.

(1) $L^{1}(\mathbf{R}^{n})$ 是没有单位元的交换 Banach 代数.

(2) $\mathfrak{M}(L^{1}(\mathbf{R}^{n})) = \mathbf{R}^{n}$.$L^{1}(\mathbf{R}^{n})$ 的 Gelfand 变换是 Fourier 变换.

(3) 令 $A(\mathbf{R}^{n}) = \{\hat{f}:\ f \in L^{1}(\mathbf{R}^{n})\}$,$\|\hat{f}\|_{A(\mathbf{R}^{n})} = \|f\|_{1}$,则 $A(\mathbf{R}^{n})$ 关于函数通常运算的和,数乘,积与范数 $\|\cdot\|_{A(\mathbf{R}^{n})}$ 成为没有单位元的交换 Banach 代数.

(4) $\mathfrak{M}(A(\mathbf{R}^{n})) = \mathbf{R}^{n}$.$A(\mathbf{R}^{n})$ 的 Gelfand 变换是恒等映射.
\end{theorem}

\begin{proof}
(1) 若 $L^{1}(\mathbf{R}^{n})$ 具有单位元 $e$,则 $e * e = e$.由此 $\hat{e} \cdot \hat{e} = \hat{e}$.因此 $\hat{e}(\xi) = 0, 1$($\xi \in \mathbf{R}^{n}$).由 Riemann-Lebesgue 引理,$\hat{e}$ 是在无穷远点取 $0$ 的连续函数,故 $\hat{e} \equiv 0$.从而由 Fourier 变换的性质,$e \equiv 0$.这与 $e$ 是单位元相矛盾.

(2) 以 $n = 1$ 证明($n > 1$ 时同理).对 $\xi \in \mathbf{R}$,令 $\chi_{\xi}(f) = \hat{f}(\xi)$,由 Fourier 变换的性质,$\chi_{\xi} \in \Delta(L^{1}(\mathbf{R}))$.又,若 $\chi_{\xi}(f) = \chi_{\xi'}(f)$($f \in L^{1}(\mathbf{R})$),则 $\xi = \xi'$,故此对应是单射.进一步,若 $\lim_{n \to \infty}\xi_{n} = \xi$,由 $\lim_{n \to \infty}\hat{f}(\xi_{n}) = \hat{f}(\xi)$,$\xi \mapsto \chi_{\xi}$ 连续.而且此对应是满射.事实上,设 $\chi \in \Delta(L^{1}(\mathbf{R}))$,则 $\chi:\ L^{1}(\mathbf{R}) \to \mathbf{C}$ 是同态.这里,$L^{1}(\mathbf{R})$ 上连续线性泛函的空间是 $L^{\infty}(\mathbf{R})$,故由 $|\chi(f)| \leq \|f\|_{1}$,存在 $\phi \in L^{\infty}(\mathbf{R})$ 使得
\begin{equation*}
\chi(f) = \int f(x)\phi(x)\frac{\dd x}{2\pi}.
\end{equation*}
进一步,$\|\phi\|_{L^{\infty}} = 1$.此时,若 $f, g \in L^{1}(\mathbf{R})$,则
\begin{align*}
\int \chi(f)g(y)\phi(y)\frac{\dd y}{2\pi} & = \chi(f)\chi(g) \\
& = \chi(f * g) \\
& = \int (f * g)(x)\phi(x)\frac{\dd x}{2\pi} \\
& = \iint g(y)f(x - y)\frac{\dd y}{2\pi}\phi(x)\frac{\dd x}{2\pi} \\
& = \int g(y)\int f(x - y)\phi(x)\frac{\dd x}{2\pi}\frac{\dd y}{2\pi} \\
& = \int g(y)\chi(f_{y})\frac{\dd y}{2\pi}
\end{align*}
成立.其中 $f_{y}(x) = f(x - y)$.因此
\begin{equation*}
\int \chi(f)\phi(y)g(y)\frac{\dd y}{2\pi} = \int \chi(f_{y})g(y)\frac{\dd y}{2\pi} \quad (g \in L^{1}(\mathbf{R}))
\end{equation*}
故得 $\chi(f)\phi(y) = \chi(f_{y})$.这里固定满足 $\chi(f) \neq 0$ 的 $f$.当 $y \to 0$ 时,$\|f_{y} - f\|_{1} \to 0$,故由 $\phi(y) = \chi(f_{y})/\chi(f)$(a.e.\ $y$),$\phi$ 可视为连续的.把 $y$ 换成 $x + y$,则
\begin{equation*}
\chi(f)\phi(x + y) = \chi(f_{x + y}) = \chi((f_{x})_{y}) = \chi(f_{x})\phi(y) = \chi(f)\phi(x)\phi(y)
\end{equation*}
故 $\phi(x + y) = \phi(x)\phi(y)$,$|\phi(x)| \leq 1$.因此 $\phi(0) = \phi(0)^{2}$.由此,$\phi(0) = 1$.从而由 $1 = \phi(0) = \phi(x)\phi(-x)$,$|\phi(x)| \leq 1$,$|\phi(x)| = 1$($x \in \mathbf{R}$).此时,存在 $\delta > 0$ 使
\begin{equation*}
\int_{0}^{\delta}\phi(x)\frac{\dd x}{2\pi} = \alpha \neq 0.
\end{equation*}
又,
\begin{equation*}
\alpha\phi(x) = \int_{0}^{\delta}\phi(y)\phi(x)\frac{\dd y}{2\pi} = \int_{0}^{\delta}\phi(x + y)\frac{\dd y}{2\pi} = \int_{x}^{x + \delta}\phi(y)\frac{\dd y}{2\pi}
\end{equation*}
故 $\phi(x)$ 可微.因此,利用
\begin{equation*}
\frac{\phi(x + h) - \phi(x)}{h} = \frac{\phi(x)\phi(h) - \phi(x)}{h} = \phi(x)\frac{\phi(h) - 1}{h}
\end{equation*}
得 $\phi'(x) = \phi'(0)\phi(x)$,$|\phi(x)| = 1$,$\phi(0) = 1$.从而存在 $\xi \in \mathbf{R}$ 使 $\phi(x) = e^{-i\xi x}$.因此
\begin{equation*}
\chi(f) = \int f(x)e^{-i\xi x}\frac{\dd x}{2\pi} = \hat{f}(\xi).
\end{equation*}
综上,$\mathbf{R} \ni \xi \mapsto \chi_{\xi} \in \Delta(L^{1}(\mathbf{R}))$ 是满射.

反之,取 $\xi_{0} \in \mathbf{R}$,$\varepsilon > 0$,设 $|\xi - \xi_{0}| < \varepsilon$.此时,取 $f \in L^{1}(\mathbf{R})$ 使 $|\hat{f}(\xi)| < |\hat{f}(\xi_{0})| = 1$($\xi \in \mathbf{R}$,$\xi \neq \xi_{0}$).这里,取适当的 $\delta > 0$ 考虑 $V = \{\chi_{\xi}:\ |\chi_{\xi}(f) - \chi_{\xi_{0}}(f)| < \delta\}$,则 $V$ 是 $\Delta(L^{1}(\mathbf{R}))$ 的开集,若 $\chi_{\xi} \in V$,可令 $|\xi - \xi_{0}| < \varepsilon$.因此 $\xi \mapsto \chi_{\xi}$ 的逆映射也连续.因此 $\mathbf{R}$ 与 $\Delta(L^{1}(\mathbf{R}))$ 同胚.从而 $f$ 的 Gelfand 变换是 $f$ 的 Fourier 变换.

(3), (4) 由 (1), (2) 容易证明.
\end{proof}

\subsection{\texorpdfstring{测度构成的交换 Banach 代数(测度代数)$M(\mathbf{T})$}{测度构成的交换 Banach 代数(测度代数)M(T)}}\label{sec:3.2.3}

这里,叙述 $\mathbf{T}$ 上有界正则 Borel 测度构成的交换 Banach 代数 $M(\mathbf{T})$ 的若干性质.$M(\mathbf{T})$ 作为交换 Banach 代数具有复杂的结构,有许多研究(\cite{ch3:2}, \cite{ch3:12}, \cite{ch3:13} 等).

定义
\begin{equation*}
L^{1}(\mathbf{T}) = \left\{f:\ \|f\|_{1}\ \bigl(= \|f\|_{L^{1}(\mathbf{T})}\bigr) = \frac{1}{2\pi}\int_{0}^{2\pi}|f(x)|\dd x < \infty\right\}.
\end{equation*}

\begin{definition}\label{def:3.2.6}

(1) 用 $M(\mathbf{T})$ 表示 $\mathbf{T}$ 上有界正则复值 Borel 测度的全体.设 $\mathcal{B}$ 为 $\mathbf{T}$ 上 Borel 集的全体,对 $\mu \in M(\mathbf{T})$,令
\begin{align*}
& |\mu|(E) = \\
& \sup\Bigl\{\sum_{j=1}^{n}|\mu(E_{j})|:\ E = \bigcup_{j=1}^{n}E_{j},\ E_{j} \in \mathcal{B},\ \{E_{j}\}_{j} \text{ 互不相交},\ n \in \mathbf{N}\Bigr\}, \\
& \|\mu\| = \|\mu\|_{M(\mathbf{T})} = |\mu|(\mathbf{T})
\end{align*}
又,当 $\alpha, \beta \in \mathbf{C}$,$\mu, \nu \in M(\mathbf{T})$ 时,定义
\begin{equation*}
\begin{aligned}
(\alpha\mu + \beta\nu)(E) &= \alpha \cdot \mu(E) + \beta \cdot \nu(E), \\
\mu * \nu(E) &= \int \mu(E - x)\dd\nu(x) \quad (E \in \mathcal{B})
\end{aligned}
\end{equation*}
则 $\alpha\mu + \beta\nu$,$\mu * \nu \in M(\mathbf{T})$.此时,$M(\mathbf{T})$ 关于范数 $\|\cdot\|$,积 $*$,和,数乘成为交换 Banach 代数.进一步,令 $\delta_{0}(E) = 1$($0 \in E$),$0$($0 \notin E$),则 $\delta_{0}$ 是 $M(\mathbf{T})$ 的单位元(称 $\delta_{0}$ 为 Dirac 测度).因此 $M(\mathbf{T})$ 是具有单位元的交换 Banach 代数.称其为 $\mathbf{T}$ 上的测度代数.

(2) 对 $\mu \in M(\mathbf{T})$,定义 Fourier-Stieltjes 变换为
\begin{equation*}
\hat{\mu}(n) = \int_{0}^{2\pi}e^{-inx}\dd\mu(x) \quad (n \in \mathbf{Z}).
\end{equation*}
\end{definition}

对 $\mu \in M(\mathbf{T})$,令 $T_{\mu}f = \int f(x)\dd\mu(x)$,则 $T_{\mu}$ 是 $C(\mathbf{T})$ 上的连续线性泛函.此时,取 $C(\mathbf{T})$ 上连续线性泛函的空间 $C(\mathbf{T})^{*}$,由 Riesz 表示定理,通过 $\mu \in M(\mathbf{T})$ 到 $T_{\mu} \in C(\mathbf{T})^{*}$ 的对应,$M(\mathbf{T})$ 与 $C(\mathbf{T})^{*}$ 作为 Banach 空间等同.以后,Riesz 表示定理视为已知(cf.\ \cite{ch3:9}, \cite{ch3:11}).

\begin{lemma}\label{lem:3.2.7}

(1) $L^{1}(\mathbf{T})$ 是 $M(\mathbf{T})$ 的闭理想.

(2) $L^{1}(\mathbf{T})$ 的极大理想空间是 $\mathbf{Z}$,且可视为 $M(\mathbf{T})$ 的极大理想空间 $\Delta(M(\mathbf{T}))$ 的开集.
\end{lemma}

\begin{proof}
(1) 对 $f \in L^{1}(\mathbf{T})$,令 $\mu_{f}(E) = \int_{E}f(x)\frac{\dd x}{2\pi}$,则 $\mu_{f} \in M(\mathbf{T})$ 且 $\|\mu_{f}\| = \|f\|_{1}$,故 $L^{1}(\mathbf{T})$ 是 $M(\mathbf{T})$ 的闭子空间.又,设 $\mathbf{T}$ 上的 Lebesgue 测度为 $dm(x) = \frac{\dd x}{2\pi}$,设 $m(E) = 0$.此时,若 $f \in L^{1}(\mathbf{T})$,$\mu \in M(\mathbf{T})$,则
\begin{equation*}
\mu_{f} * \mu(E) = \int_{0}^{2\pi}\mu_{f}(E - x)\dd\mu(x) = 0
\end{equation*}
故得 $f * \mu = \mu_{f} * \mu \in L^{1}(\mathbf{T})$.因此 $L^{1}(\mathbf{T})$ 是 $M(\mathbf{T})$ 的理想.

(2) $L^{1}(\mathbf{T})$ 的极大理想空间可与定理~\ref{thm:3.2.5} 同样地证明与 $\mathbf{Z}$ 等同(利用 $\phi$ 具有 $2\pi$ 周期性).因此 $\Delta(L^{1}(\mathbf{T})) = \mathbf{Z}$.由 $L^{1}(\mathbf{T})$ 是 $M(\mathbf{T})$ 的理想,在 $\Delta(M(\mathbf{T}))$ 中 $\mathbf{Z}$ 是开集.
\end{proof}

\begin{theorem}[Williamson \cite{ch3:55}]\label{thm:3.2.8}
存在 $\mu \in M(\mathbf{T})$ 使 $\Sp(\mu, M(\mathbf{T}))$ 真包含 $\hat{\mu}(\mathbf{Z})$ 的闭包.由此,$\overline{\mathbf{Z}}$ 是 $\Delta(M(\mathbf{T}))$ 的真子集.
\end{theorem}

\begin{proof}
由后述的注意~\ref{note:3.3.12},存在 $\mu \in M(\mathbf{T})$ 使得 $\mu(\mathbf{T}) = 1$,$\mu \geq 0$(称这样的 $\mu$ 为概率测度),$\hat{\mu}$ 取实数值,且任意一点的测度为 $0$.进一步,此测度可取为:当令 $\mu^{n} = \mu^{n-1} * \mu$($n > 1$)时,$\delta_{0}, \mu, \mu^{2}, \mu^{3}, \ldots$ 具有落在互不相交的集合上的测度.此时,对 $\lambda \in \mathbf{C}$,$|\lambda| = 1$,
\begin{align*}
& \|(\lambda\delta_{0} + \mu)^{n}\| = \left\|\sum_{k=0}^{n}\binom{n}{k}\lambda^{k}\mu^{n-k}\right\| \\
& \qquad = \sum_{k=0}^{n}\binom{n}{k}|\lambda|^{k}\|\mu\|^{n-k} \\
& \qquad = \sum_{k=0}^{n}\binom{n}{k} = 2^{n}
\end{align*}
成立.因此 $\lim_{n \to \infty}\|(\lambda\delta_{0} + \mu)^{n}\|^{\frac{1}{n}} = 2$,故存在 $\chi \in \Delta(M(\mathbf{T}))$ 使 $|\chi(\lambda\delta_{0} + \mu)| = 2$.从而得 $\chi(\mu) = \lambda$.因此,按 $\mu$ 的取法,$\overline{\hat{\mu}(\mathbf{Z})} \subsetneq \Sp(\mu, M(\mathbf{T}))$.由此与补题~\ref{lem:3.2.7} 得 $\overline{\mathbf{Z}} \subsetneq \Delta(M(\mathbf{T}))$.
\end{proof}

下面,用定理~\ref{thm:3.2.8} 的 $\mu$,令 $\nu = \mu^{2} + \delta_{0}$,即可容易得到.

\begin{corollary}\label{cor:3.2.9}
存在满足 $\hat{\nu} \geq 1$ 的 $\nu \in M(\mathbf{T})$,且 $1/\hat{\nu}$ 不是 Fourier-Stieltjes 变换.
\end{corollary}

这些结果本质上是由 Wiener-Pitt \cite{ch3:54} 于 1938 年得到的,把这种现像称为 Wiener-Pitt 现象.在这个方向上,F.\ Parreau 得到了更精密的结果,是对下述 Katznelson 问题 \cite{ch3:6} 的否定解决.

\textbf{Katznelson 问题}．对 $\mu \in M(\mathbf{T})$,若 $\Sp(\mu, M(\mathbf{T}))$ 是实数的子集,则 $\hat{\mu}(\mathbf{Z})$ 在 $\Sp(\mu, M(\mathbf{T}))$ 中稠密吗?

\begin{theorem}[Parreau \cite{ch3:36}]\label{thm:3.2.10}
存在 $\mu \in M(\mathbf{T})$ 满足:
\begin{equation*}
\mathbf{R} \supset \Sp(\mu, M(\mathbf{T})) \supsetneqq \overline{\hat{\mu}(\mathbf{Z})}.
\end{equation*}
\end{theorem}

\subsection{\texorpdfstring{交换 Banach 代数 $\mathfrak{M}_p(\mathbf{T})$}{交换 Banach 代数 M_p(T)}}\label{sec:3.2.4}

设 $1 \leq p < \infty$.令 $L^{p}(\mathbf{T}) = \{f:\ \|f\|_{p} = \left(\frac{1}{2\pi}\int_{0}^{2\pi}|f(x)|^{p}\dd x\right)^{1/p} < \infty\}$,则 $L^{p}(\mathbf{T})$ 由通常的运算成为 Banach 空间.对 $\mathbf{T}$ 上的函数 $f(x)$,令 $\tau_{y}f(x) = f(x - y)$,称满足 $T\tau_{y} = \tau_{y}T$($y \in \mathbf{T}$)的 $L^{p}(\mathbf{T})$ 上的连续线性算子 $T$ 为乘子算子,其全体记为 $\mathfrak{M}_{p}(\mathbf{T})$.此时,$\mathfrak{M}_{p}(\mathbf{T})$ 以算子范数 $\|T\|_{\mathfrak{M}_{p}(\mathbf{T})} = \sup\{\|Tf\|_{p}:\ \|f\|_{p} \leq 1\}$ 与自然的运算成为 Banach 空间.又,利用对偶性可证明 $T(f * g) = Tf * g$($f, g \in C(\mathbf{T})$).进一步,具有下列性质(cf.\ \cite{ch3:9}).

\begin{lemma}\label{lem:3.2.11}
设 $1 \leq p \leq 2$.

(1) 当 $T \in \mathfrak{M}_{p}(\mathbf{T})$ 时,存在 $\hat{T} \in \ell^{\infty}$ 使 $(Tf)^{\wedge}(n) = \hat{T}(n)\hat{f}(n)$($f \in C(\mathbf{T})$).这里,$\ell^{\infty}$ 表示由 $\mathbf{Z}$ 上的有界数列 $\{a_{n}\}$ 构成的,以 $\|\{a_{n}\}\|_{\ell^{\infty}} = \sup_{n}|a_{n}|$ 为范数的 Banach 空间.又,记 $M_{p}(\mathbf{Z}) = \{\hat{T}:\ T \in \mathfrak{M}_{p}(\mathbf{T})\}$,称 $M_{p}(\mathbf{Z})$ 的元为乘子.$M_{p}(\mathbf{Z})$ 的范数由 $\mathfrak{M}_{p}(\mathbf{T})$ 自然引入.

(2) $\mathfrak{M}_{1}(\mathbf{T})$ 作为 Banach 空间与 $M(\mathbf{T})$ 等同.

(3) $M_{2}(\mathbf{Z})$ 作为 Banach 空间与 $\ell^{\infty}$ 等同.
\end{lemma}

\begin{proof}
(1) 当 $f(x) = e^{-inx}$ 时,令 $\hat{T}(n) = \int_{0}^{2\pi}Tf(x)e^{-inx}\frac{\dd x}{2\pi}$,由 Hölder 不等式,$|\hat{T}(n)| \leq \|T\|_{\mathfrak{M}_{p}}$.此时,$T(e^{-inx}) = \hat{T}(n)e^{inx}$,故对 $f \in C(\mathbf{T})$,
\begin{equation*}
\widehat{Tf}(n) = \hat{T}(n)\hat{f}(n) \quad (n \in \mathbf{Z}).
\end{equation*}
又,由三角多项式全体在 $L^{p}(\mathbf{T})$ 中稠密,$\{\hat{T}(n)\}$ 唯一确定.

(2) 若 $\mu \in M(\mathbf{T})$,由 $T_{\mu}f = \mu * f$,$T_{\mu}f \in L^{1}(\mathbf{T})$($f \in L^{1}(\mathbf{T})$).又,由 Fubini 定理,$\|T_{\mu}f\|_{1} = \int\left|\int f(x - y)\dd\mu(y)\right|\frac{\dd x}{2\pi} \leq \|\mu\|\|f\|_{1}$,得 $\|T_{\mu}\|_{\mathfrak{M}_{1}} \leq \|\mu\|$.反之,当 $T \in \mathfrak{M}_{1}(\mathbf{T})$ 时,对 $\mathbf{T}$ 上的三角多项式 $g$,对应 $\sum_{n}\hat{g}(n)\hat{T}(n) = \int T(e^{-inx})g(-x)\frac{\dd x}{2\pi}$,由 (1),$\left|\sum_{n}\hat{g}(n)\hat{T}(n)\right| \leq \|T\|_{\mathfrak{M}_{1}}\|g\|_{\infty}$,故由 Riesz 表示定理,存在 $\mu \in M(\mathbf{T})$ 使 $\int g(-x)\dd\mu(x) = \sum_{n}\hat{g}(n)\hat{T}(n)$.特别地,当 $g(x) = e^{inx}$ 时,$\hat{\mu}(n) = \hat{T}(n)$,故可写为 $T = T_{\mu}$.因此 $\|\mu\| \leq \|T\|_{\mathfrak{M}_{1}}$.综上,$\|\mu\| = \|T_{\mu}\|_{\mathfrak{M}_{1}}$,$\mathfrak{M}_{1}(\mathbf{T})$ 作为 Banach 空间与 $M(\mathbf{T})$ 等同.

(3) 若 $\{a_{n}\} \in \ell^{\infty}$,则当 $f \in L^{2}(\mathbf{T})$ 时,由 Parseval 等式,
\begin{equation*}
\sum\left|a_{n}\hat{f}(n)\right|^{2} \leq \|\{a_{n}\}\|_{\ell^{\infty}}^{2}\sum|\hat{f}(n)|^{2} = \|\{a_{n}\}\|_{\ell^{\infty}}^{2}\|f\|_{2}^{2} < \infty
\end{equation*}
故存在 $g \in L^{2}(\mathbf{T})$ 使 $\hat{g}(n) = a_{n}\hat{f}(n)$($n \in \mathbf{Z}$).因此,由 $\widehat{Tf}(n) = a_{n}\hat{f}(n)$ 确定 $T \in \mathfrak{M}_{2}(\mathbf{T})$,且可知 $\|T\|_{\mathfrak{M}_{2}} \leq \|\{a_{n}\}\|_{\ell^{\infty}}$.反之,若 $T \in \mathfrak{M}_{2}(\mathbf{T})$,像 (2) 那样令 $\hat{T}(n) = \int T(e^{-inx})e^{-inx}\frac{\dd x}{2\pi}$,由 Schwarz 不等式可证明 $\hat{T} \in \ell^{\infty}$ 且 $\|\{\hat{T}(n)\}\|_{\ell^{\infty}} \leq \|T\|_{\mathfrak{M}_{2}}$.
\end{proof}

下面的证明由补题~\ref{lem:3.2.11} 容易得到.

\begin{theorem}\label{thm:3.2.12}
设 $1 \leq p \leq 2$.$\mathfrak{M}_{p}(\mathbf{T})$ 关于算子的复合,成为以 $I = \delta_{0}$ 为单位元的交换 Banach 代数.
\end{theorem}

\subsection{奇异积分算子构成的交换 Banach 代数}\label{sec:3.2.5}

这里,对 $n \geq 2$,只介绍奇异积分算子全体构成交换 Banach 代数这一 Calderón-Zygmund \cite{ch3:14} 的结果.证明等细节请参照 \cite{ch3:14}.

\begin{definition}\label{def:3.2.13}
$\Omega(x)$ 是 $\mathbf{R}^{n} \setminus \{0\}$ 上的函数,设 $S^{n-1}$ 为 $\mathbf{R}^{n}$ 的单位球面,对 $1 < q < \infty$ 满足下列性质:

(1) $\Omega(rx) = \Omega(x)$($x \in S^{n-1}$,$r > 0$)

(2) $\Omega \in L^{q}(S^{n-1})$

(3) $\int_{S^{n-1}}\Omega(x')\dd\sigma(x') = 0$.其中 $\dd\sigma$ 表示 $S^{n-1}$ 上的面积元素,设 $L^{q}(S^{n-1}) = \{f:\ \|f\|_{L^{q}(S^{n-1})} = \left(\int_{S^{n-1}}|f(x')|^{q}\dd\sigma(x')\right)^{1/q} < \infty\}$.
\end{definition}

令 $K(x) = \Omega(x)/|x|^{n}$($x \neq 0$)时,奇异积分算子
\begin{equation*}
K * f(x) = \lim_{\varepsilon \to 0}\int_{|x - y| > \varepsilon}K(x - y)f(y)\dd y
\end{equation*}
是 $L^{p}(\mathbf{R}^{n})$($1 < p < \infty$)上的有界线性算子,其 Fourier 变换为
\begin{equation*}
\mathcal{F}(K * f) = \mathcal{F}(K)\mathcal{F}(f) \quad \left(f \in C_{c}^{\infty}(\mathbf{R}^{n})\right),
\end{equation*}
且已知在几乎所有点上 $|\mathcal{F}(K)| \leq C\|K\|_{L^{q}(S^{n-1})}$ 成立(cf.\ \cite{ch3:5}).其中 $C_{c}^{\infty}(\mathbf{R}^{n})$ 是具有紧支集的无限次连续可偏导函数的全体.把上述 $K$ 的全体记为 $\mathcal{K}$.此时,对 $1 < q < \infty$,定义下列奇异积分算子的族
\begin{equation*}
A_{q} = \{T:\ Tf = \alpha f + K * f,\ \alpha \in \mathbf{C},\ K \in \mathcal{K}\}.
\end{equation*}
进一步,对 $T \in A_{q}$,$Tf = \alpha f + K * f$,令
\begin{equation*}
|T|_{q} = |\alpha| + \|K\|_{L^{q}(S^{n-1})}.
\end{equation*}
此时,下式成立:

\begin{theorem}[Calderón-Zygmund]\label{thm:3.2.14}
$A_{q}$ 关于算子的通常运算与范数 $|\cdot|_{q}$,成为具有单位元的交换 Banach 代数.
\end{theorem}

又,关于 $A_{q}$ 的极大理想空间,已知下列结论.

\begin{theorem}\label{thm:3.2.15}
具有单位元的交换 Banach 代数 $A_{q}$ 的极大理想空间 $\mathfrak{M}(A_{q})$ 与 $S^{n-1}$ 等同.
\end{theorem}
